\section{Conclusiones y trabajo futuro}

Esta investigación ha demostrado que las arquitecturas de detección de objetos basadas en Transformers, específicamente RF-DETR, representan una solución superior para la monitorización automatizada de fauna en manadas densas, superando las limitaciones inherentes a los métodos tradicionales basados en CNN y NMS. Los hallazgos clave se resumen a continuación:

\begin{itemize}
    \item \textbf{Superioridad de la Arquitectura sin NMS:} La predicción directa de conjuntos (\textit{end-to-end}) de RF-DETR resolvió eficazmente el problema de subconteo por oclusión. El modelo óptimo, \textbf{RF-DETR Small}, alcanzó un \textbf{F1-Score global de 90.24\%}, superando al baseline HerdNet (83.26\%) y reduciendo el Error Medio Absoluto (MAE) de conteo en un \textbf{39.47\%}.

    \item \textbf{Eficiencia Óptima del Modelo Small:} Se identificó un punto de saturación en el rendimiento donde aumentar la complejidad del modelo no se traduce en ganancias significativas. La variante \textbf{Small} (32M parámetros) superó ligeramente a la variante \textbf{Large} (128M parámetros) en precisión global, siendo \textbf{56\% más rápida} en inferencia media (193 ms vs 442 ms), consolidándose como la opción ideal para el despliegue.

    \item \textbf{Validación del Hard Negative Mining:} El entrenamiento en dos fases fue crítico. La minería de negativos difíciles (HNM) transformó a RF-DETR Small de un modelo con precisión inaceptable (26.15\%) a la solución más precisa (93.85\%), demostrando la plasticidad del modelo para discriminar fondos complejos sin aumentar su tamaño.

    \item \textbf{Mejora en Especies Minoritarias:} El uso de \textit{backbones} visuales robustos (DINOv2) mejoró significativamente la detección de clases subrepresentadas. En el caso crítico del Cobo de agua (2.4\% del dataset), RF-DETR Small mantuvo un F1-Score robusto superior al 70\%, validando su capacidad de generalización con pocos datos.

    \item \textbf{Limitación de la Resolución Mínima:} La variante Nano demostró ser insuficiente para esta tarea, colapsando en la detección de especies clave como el Elefante, lo que establece una resolución de entrada mínima de $512 \times 512$ como requisito para capturar patrones complejos en imágenes aéreas.\newline
\end{itemize}

\subsubsection*{Para el trabajo futuro, se proponen dos direcciones principales}

\begin{enumerate}
    \item \textbf{Mitigación Avanzada del Desbalance de Clases:}
    Aunque RF-DETR mejoró la detección base, persiste un margen de mejora en las clases minoritarias. Se propone incorporar estrategias de ponderación de pérdida (\textit{Weighted Loss}) en la función de costo, replicando la estrategia de asignación agresiva de pesos que demostró ser efectiva en el entrenamiento de HerdNet. Alternativamente, se sugiere explorar técnicas de sobre-muestreo (\textit{Oversampling}) para equilibrar la exposición del modelo a estas especies durante el entrenamiento.

    \item \textbf{Optimización de Inferencia mediante cuantización:}
    Para facilitar el despliegue en dispositivos de borde, es necesario explorar la cuantización del modelo a INT8. Dado que las arquitecturas basadas en Transformers pueden sufrir degradación de rendimiento con una cuantización ingenua \cite{rf-detr}, se recomienda implementar esquemas de \textit{Post-Training Quantization} (PTQ) calibrados \cite{fima-q} o \textit{Quantization-Aware Training} (QAT). El objetivo es reducir la huella de memoria sin que la pérdida de precisión numérica afecte desproporcionadamente la localización en manadas densas \cite{hqod}.

    \item \textbf{Entrenamiento con DINOV3 como extractor de características:}

    Si bien la implementación actual de RF-DETR utiliza DINOv2 \cite{dinov2}, la evolución reciente de la arquitectura sugiere la adopción de DINOv3 \cite{dinov3}. Se recomienda evaluar la migración hacia este nuevo backbone, dado que ha demostrado superar a su predecesor en tareas de detección, logrando un rendimiento superior con una menor cantidad de parámetros y mayor eficiencia computacional \cite{deimv2}
\end{enumerate}
